% !TEX program = xelatex
\documentclass{ctexrep}

\usepackage{fontawesome}
\usepackage{ecnu}
\author{张yadong}
\date{\today}
\title{人文社科类学位论文写作—— \LaTeX 轻量级应用}
\begin{document}
\maketitle
\tableofcontents



\section{写作缘由}
\label{sec:orgddf61dc}
\LaTeX{}的论文排版无可替代，而学校不会专门教排版设计。
针对人文社科类学位论文写作，从基本的文字编辑，数据分析，进度管理，版本控制到最后的排版设计，笔者写下这篇文章，
希望能够帮助大家完成自己的
毕业论文。
笔者默认你使用\faWindows~ Windows10操作系统，\faApple~ Mac 原理一样，Linux\ldots{}\faLinux 嗯，你应该比我知道的多。
\section{文字编辑}
\label{sec:org63c45b2}
\subsection{Org-mode 安装}
\label{sec:orgaaef9fc}
注：下列键盘符号 C 代表 Ctrl 键， M 代表 Alt 键，RET 代表回车，小写字母就是实际字母。

\begin{itemize}
\item 下载\href{https://www.gnu.org/software/emacs/download.html}{Emacs} 26.2，安装最新版 \href{https://orgmode.org/elpa.html}{Org-mode} 9.2.3，官网说的配置文件 \textbf{init file} 需要新建：
\item 打开 \textbf{runemaca.exe} ，C-x C-f 输入 \textbf{.emacs} ，新建完成
\item 按照官网介绍操作，安装完成后，M-x org-version，应该就是与官网一致的版本。
\end{itemize}
\subsection{基本操作}
\label{sec:orgad93b47}
参照国外大佬的\href{https://pan.baidu.com/s/1p6CRrnt6c0WrROvLW0BjRA }{视频（提取码：26qz ）}，主要记 \textbf{快捷键} （输入字母可以识别，但不是最新的，比如大写字母都改小写了），介绍顺序在\href{https://orgmode.org/worg/org-tutorials/org-screencasts/org-mode-google-tech-talk.html}{这儿}。
\subsection{特殊写作环境}
\label{sec:orgf18cb22}
写作前，开头加入这段文本:
\begin{verbatim}
# -*- coding: utf-8 -*-
\end{verbatim}

保证你的文字导出时不会乱码。

1.摘要

\begin{verbatim}
#+begin_abstract
摘要测试
#+end_abstract
\end{verbatim}

2.代码

C-c C-,

默认为verbatim环境，
如果要用 minted 宏包，确认你有安装了Python包pygments，建议安装 \textbf{Anaconda} ，然后在 \#+begin
前一行加上

\begin{verbatim}
#+ATTR_LATEX: :options org-latex-minted-options
\end{verbatim}

\section{数据分析}
\label{sec:org40f6abb}
SPSS良心14天试用期够用了，如果你经常要写论文，推荐学习一下 R 语言。
\section{{\bfseries\sffamily TODO} 使用 Org-mode 的GTD工作流程}
\label{sec:org65c1eca}
\section{版本控制（可选）}
\label{sec:org567cdd4}
\subsection{安装 \href{https://git-scm.com/downloads}{Git} 2.21}
\label{sec:org8cd1a95}
\subsection{\href{https://www.liaoxuefeng.com/wiki/896043488029600}{廖雪峰的Git教程}}
\label{sec:org250cbff}
\section{协同写作}
\label{sec:org6b80c8e}
C-c C-e 导出utf-8文本，供导师修改（反正排版是最后做的）。
\section{排版设计}
\label{sec:orgf6e0dca}
笔者在设计过程中发现Org-mode中用 Xe\LaTeX 写中文文档很早就有人写配置（添加在 \textbf{.emacs} 文件的）了，
现状笔者添加到 ox-latex.el里，然后编译：M-x byte-compile-file，重新打开 runemacs.exe （推荐
添加到桌面快捷方式）就可以使用了。
所有文本写完之后，跳到开头，C-c C-e \#，输入latex，应该可以看到一下信息：

\begin{verbatim}
#+latex_class: ctexrep
#+latex_class_options:<默认为[12pt, a4paper],可以自己设置>
#+latex_header:{\usepackage{学校 \LaTeX 模板样式（.sty结尾），我用的 ecnu}
#+latex_header_extra:
#+description:
#+keywords:
#+subtitle:
#+author: 
#+title: <默认显示你的文件名>
#+latex_compiler: xelatex
#+date: \today 
\end{verbatim}

author，title，date可以删掉，如果有学校封面。

表格，图片一律使用图片导入：C-c C-l file，选择你的图片路径即可。
示例：
\begin{center}
\includegraphics[width=.9\linewidth]{org-mode-unicorn-logo.png}
\end{center}
\subsection{安装 \href{https://zhuanlan.zhihu.com/p/64555335}{\TeX{} Live}（强烈推荐）2019}
\label{sec:orgfd45bc8}
安装完成后，打开你的 \textbf{.org} 文件，C-c C-e l o，即可打开编译好的 \textbf{pdf} 文件，当然可以自己改一下
\textbf{.tex} 文件，这样排版工作量就会少很多。
\section{参考链接}
\label{sec:orgb91b710}
\subsection{{\bfseries\sffamily TODO} 基于 gbt-7714-2015 格式参考文献编译}
\label{sec:org4da2f8a}

[1] \url{https://www.reddit.com/r/emacs/comments/4k1lp2/noob\_question\_how\_to\_set\_locales\_and\_encoding\_for/}

[2] \url{https://www.cnblogs.com/wangkangluo1/archive/2012/02/04/2337705.html}

[3] \url{http://www.cnblogs.com/visayafan/archive/2012/06/16/2552023.html}

[4] \url{https://xiaoguo.net/wiki/org-mode-book.html}

[5] \url{https://orgmode.org/manual/index.html\#SEC\_Contents}

[6] \url{https://orgmode.org/worg/org-tutorials/org-screencasts/org-mode-google-tech-talk.html\#sec-2}

[7] \url{https://zhuanlan.zhihu.com/p/64555335}
\end{document}
